\documentclass{article}
\linespread{1.3}
\usepackage[margin=50pt]{geometry}
\usepackage{amsmath, amsthm, amssymb, amsthm, tikz, fancyhdr, float, graphicx, spverbatim, systeme}
\pagestyle{fancy}
\renewcommand{\headrulewidth}{0pt}
\newcommand{\changefont}{\fontsize{15}{15}\selectfont}

\fancypagestyle{firstpageheader}{
  \fancyhead[R]{
    \changefont
    \parbox[t]{4cm}{ % Adjust width as needed
      Michael Huang\\
      EN.555.644.81\\
      Homework N
    }
  }
}

\begin{document}

\thispagestyle{firstpageheader}
{\Large 

\section*{11.12}
The price of a non-dividend paying stock is \$19 and the price of a three-month European call option on the stock with a strike price of \$20 is \$1. The risk-free rate is 4\% per annum. What
is the price of a three-month European put option with a strike price of \$20?


}
{\Large 

\section*{11.19}
The price of a European call that expires in six months and has a strike price of \$30 is \$2.
The underlying stock price is \$29, and a dividend of \$0.50 is expected in two months and
again in five months. Interest rates (all maturities) are 10\%. What is the price of a European
put option that expires in six months and has a strike price of \$30?


}
{\Large 

\section*{11.20}
Explain the arbitrage opportunities in Problem 11.19 if the European put price is \$3.

}
{\Large 

\section*{11.23}
Prove the result in equation (11.7). (Hint: For the first part of the relationship consider (a) a
portfolio consisting of a European call plus an amount of cash equal toK and (b) a portfolio
consisting of an American put option plus one share.)


}
{\Large 

\section*{11.24}
Prove the result in equation (11.11). (Hint: For the first part of the relationship consider (a)
a portfolio consisting of a European call plus an amount of cash equal toD K+ and (b) a
portfolio consisting of an American put option plus one share.)


}
{\Large 

\section*{11.26}
Use the software DerivaGem to verify that Figures 11.1 and 11.2 are correct. 


}
{\Large 

\section*{11.31}
Suppose that \(c_1, c_2,\) and \(c_3\) are the prices of European call options with strike prices \(K_1, K_2,\) and \(K_3\), respectively, where \(K_3 > K_2 > K_1\) and \(K_3 - K_2 = K_2 - K_1\). All options have the same maturity. Show that 
\[
  c_2 \leq 0.5(c_1 + c_3)
\]
(Hint: Consider a portfolio that is long one option with strike price \(K_1\), long one option with strike price \(K_3\), and short two optiosn with strike price \(K_2\).)


}
{\Large 

\section*{.}
Consider an option on a stock when the stock price is \$41, the strike price of \$40, the risk-free rate is 6\%, the volatility is 35\%,and the time to maturity is 1 year. Assume that a dividend of \$0.50 is expected after six months. 

\subsection*{(a)} 
Use DerivaGem to value the option assuming it is a European call.

\subsection*{(b)}
Use DerivaGem to value the option assuming it is a European put.

\subsection*{(c)}
Verify that put-call parity holds.

\subsection*{(d)}
Explore using DerivaGem what happens to the price of the options as the time to maturity becomes very large and there are no dividends. Explain your results. 

}

\end{document}